% Estilo híbrido académico moderno (tipo O'Reilly):
% Header: título del capítulo (par) / título de sección (impar)
% Footer: número de página alineado al borde exterior
\usepackage{fancyhdr}
\pagestyle{fancy}
\fancyhf{} % limpia todos los campos
\fancyhead[LE]{\nouppercase{\small\itshape\leftmark}}
\fancyhead[RO]{\nouppercase{\small\itshape\rightmark}}
\fancyfoot[LE,RO]{\thepage}
\renewcommand{\headrulewidth}{0.4pt}
\renewcommand{\footrulewidth}{0pt}

% Páginas plain (inicio de capítulo): solo número de página al pie exterior
\fancypagestyle{plain}{%
  \fancyhf{}%
  \fancyfoot[LE,RO]{\thepage}%
  \renewcommand{\headrulewidth}{0pt}%
  \renewcommand{\footrulewidth}{0pt}%
}

% Páginas en blanco automáticas — leyenda según el idioma del documento
% (\languagename: babel/polyglossia; español → ES, cualquier otro → EN)
\usepackage{etoolbox}
% Páginas en blanco automáticas (forzadas por book class para que cada
% capítulo inicie en recto): leyenda centrada en gris discreto que aclara
% al lector que la página no es un error de impresión.
\makeatletter
\renewcommand*{\cleardoublepage}{%
  \clearpage
  \if@twoside
    \ifodd\c@page\else
      \thispagestyle{empty}%
      \vspace*{\fill}%
      \begin{center}\textcolor{gray!60}{\small\itshape
        \ifdefstring{\languagename}{spanish}%
          {Esta página se dejó intencionalmente en blanco.}%
          {This page intentionally left blank.}}\end{center}%
      \vspace*{\fill}%
      \newpage
    \fi
  \fi
}
\makeatother

% --- Fallbacks para caracteres Unicode no presentes en STIX Two Text ---
% Mapean glifos (≠ → ✓ △ etc.) a sus equivalentes math, que sí están en
% STIX Two Math. Sin esto, xelatex emite warnings y los caracteres
% aparecen como cuadrados vacíos en el PDF.
\usepackage{newunicodechar}
\newunicodechar{≠}{\ensuremath{\neq}}
\newunicodechar{→}{\ensuremath{\rightarrow}}
\newunicodechar{←}{\ensuremath{\leftarrow}}
\newunicodechar{↔}{\ensuremath{\leftrightarrow}}
\newunicodechar{⇄}{\ensuremath{\leftrightarrow}}
\newunicodechar{⟺}{\ensuremath{\Longleftrightarrow}}
\newunicodechar{⇔}{\ensuremath{\Leftrightarrow}}
\newunicodechar{★}{\ensuremath{\bigstar}}
\newunicodechar{✓}{\ensuremath{\checkmark}}
\newunicodechar{✗}{\ensuremath{\times}}
\newunicodechar{∞}{\ensuremath{\infty}}
\newunicodechar{≈}{\ensuremath{\approx}}
\newunicodechar{≤}{\ensuremath{\leq}}
\newunicodechar{≥}{\ensuremath{\geq}}
\newunicodechar{△}{\ensuremath{\triangle}}
\newunicodechar{▲}{\ensuremath{\blacktriangle}}
\newunicodechar{▼}{\ensuremath{\blacktriangledown}}
\newunicodechar{▽}{\ensuremath{\bigtriangledown}}
\newunicodechar{◇}{\ensuremath{\diamond}}
\newunicodechar{◆}{\ensuremath{\blacklozenge}}
